\documentclass[11pt,a4paper]{article}

\usepackage[a4paper,margin=2.5cm]{geometry}
\usepackage[french]{babel}
\usepackage{amsmath,amssymb,amsfonts,mathtools}
\usepackage{graphicx}
\usepackage{xcolor}
\usepackage{booktabs}
\usepackage{hyperref}
\usepackage{physics}
\usepackage{bm}
\usepackage{microtype}

\hypersetup{
    colorlinks=true,
    linkcolor=blue!60!black,
    citecolor=blue!60!black,
    urlcolor=blue!60!black
}

\title{
\textbf{Ondes Gravitationnelles dans la Gravité Quantique Bottom-Up}\\
\large Modes Informationnels du Graphe d'Intrication
}

\author{Farid Hamdad}
\date{\today}

\begin{document}

\maketitle

\begin{abstract}
En Gravité Quantique Bottom-Up (BuP), l'espace-temps n'est pas supposé fondamental.
Le point de départ est un état quantique \(|\Psi\rangle\), à partir duquel on calcule
les matrices densité réduites, les entropies et les informations mutuelles. Le
graphe d'intrication
\[
W_{ij} = I(i:j)
\]
n'est donc pas une structure de fond, mais un objet information-géométrique dérivé.
Cet article étudie les ondes gravitationnelles dans BuP. Contrairement à l'image
standard où la métrique est première, une onde gravitationnelle BuP est décrite
comme une perturbation collective du graphe d'intrication,
\[
\delta W_{ij}(t)
\;\rightarrow\;
\delta d_{ij}^{\rm ent}(t)
\;\rightarrow\;
\delta g_{\mu\nu}^{\rm eff}(t).
\]
Nous présentons une séquence de prototypes numériques et variationnels.
Premièrement, des diagnostics à jauge fixée retrouvent les deux polarisations
quadripolaires \(+\) et \(\times\). Deuxièmement, des perturbations propagatives
imposées se propagent avec la vitesse de phase attendue. Troisièmement, des
sources dynamiques localisées émettent des ondes de graphe avec
\(v_+ \simeq v_\times \simeq 1\) en unités de graphe. Enfin, la linéarisation d'une
action BuP effective autour d'un graphe d'équilibre donne des modes propres hessiens
satisfaisant
\[
\omega^2 \simeq c_{\rm graphe}^2 \, k^2 + m_{\rm eff}^2,
\]
avec des ajustements de dispersion \(R^2 \simeq 0{,}999\). La vitesse de propagation
est contrôlée par la rigidité informationnelle du graphe et peut être calibrée
sur la vitesse physique de la lumière dans la limite einsteinienne lisse.

Nous testons également une extension cosmologique exploratoire dans laquelle
les ondes gravitationnelles primordiales traversent une phase fortement
intriquée avec \(c_{\rm GW} < c\). Le modèle phénoménologique suggère une bande PTA
inchangée mais une amplification dans la bande LISA, maximale pour
\(\alpha_{\rm early} \simeq 0{,}35\), où accumulation temporelle et amortissement
informationnel s'équilibrent.

Ces résultats soutiennent l'interprétation des ondes gravitationnelles comme
modes informationnels du réseau d'intrication.
\end{abstract}

\section{Introduction}

En relativité générale, les ondes gravitationnelles sont des perturbations de la
métrique,
\[
g_{\mu\nu} \;\rightarrow\; g_{\mu\nu} + h_{\mu\nu}.
\]

En Gravité Quantique Bottom-Up, la métrique n'est pas fondamentale. Elle émerge
de la structure informationnelle d'un état quantique. L'objet fondamental est un
état \(|\Psi\rangle\). À partir de cet état, on calcule les matrices densité
réduites \(\rho_i\), \(\rho_{ij}\), les entropies \(S_i\), \(S_{ij}\), et les
informations mutuelles
\[
W_{ij} = I(i:j) = S_i + S_j - S_{ij}.
\]

Le graphe d'intrication \(W_{ij}\) définit une structure information-géométrique.
Les distances sont reconstruites à travers une distance d'intrication,
\[
d_{ij}^{\rm ent} = -\ell_0 \log\left(\frac{W_{ij}}{W_0}\right),
\]
et une métrique effective est reconstruite dans la limite lisse,
\[
W_{ij}
\;\rightarrow\;
d_{ij}^{\rm ent}
\;\rightarrow\;
g_{\mu\nu}^{\rm eff}.
\]

Une onde gravitationnelle BuP n'est donc pas primitivement une ondulation d'une
métrique préexistante. C'est une perturbation collective du graphe d'intrication :
\[
\delta W_{ij}(t)
\;\rightarrow\;
\delta d_{ij}^{\rm ent}(t)
\;\rightarrow\;
\delta g_{\mu\nu}^{\rm eff}(t).
\]

Cet article teste cette idée en utilisant quatre prototypes successifs :
\begin{enumerate}
    \item extraction à jauge fixée des deux polarisations ;
    \item ondes de graphe propagatives imposées ;
    \item sources dynamiques localisées ;
    \item modes hessiens d'une action BuP effective.
\end{enumerate}

\section{Ondes gravitationnelles BuP}

La perturbation métrique linéarisée standard s'écrit
\[
h_{\mu\nu} = \delta g_{\mu\nu}.
\]

Dans BuP, cet objet est la projection géométrique finale d'une perturbation
informationnelle plus profonde :
\[
\delta W_{ij}
\;\rightarrow\;
\delta d_{ij}^{\rm ent}
\;\rightarrow\;
h_{\mu\nu}^{\rm eff}.
\]

Ceci motive la définition d'une onde gravitationnelle BuP comme un mode collectif
propagatif du graphe d'intrication. Dans cette perspective, le graviton n'est pas
une particule ponctuelle fondamentale. C'est le quantum d'un mode collectif du
réseau d'intrication.

\section{Diagnostics de polarisation}

Dans le premier prototype, la perturbation du graphe d'intrication est projetée
sur des composantes quadripolaires à jauge fixée,
\[
q_+(t), \qquad q_\times(t).
\]

Les modes plus et croix sont reconstruits à partir de la dépendance angulaire
des perturbations relatives de distance d'intrication :
\[
\frac{\delta d_{ij}^2}{d_{ij}^2}
\simeq
q_+(t)\cos(2\theta_{ij})
+
q_\times(t)\sin(2\theta_{ij}).
\]

L'analyse à jauge fixée sépare les deux polarisations et confirme que les
perturbations du graphe BuP peuvent reproduire la structure quadripolaire
habituelle des ondes gravitationnelles.

\section{Ondes de graphe propagatives}

Le second prototype impose une perturbation propagative du graphe d'intrication,
\[
\delta W_{ij}(x,t) \sim A \cos(kx - \omega t).
\]

La phase quadripolaire résultante est mesurée sur des tranches spatiales. La
pente de phase reconstruit le nombre d'onde imposé :
\[
k_{\rm fit} \simeq k_{\rm input}.
\]

Ceci vérifie que le graphe d'intrication supporte des modes quadripolaires
propagatifs.

\section{Émission par source dynamique}

Le troisième prototype remplace l'onde propagative imposée par une source
localisée oscillante. La réponse est mesurée sur des anneaux concentriques.

Pour la polarisation plus, l'ajustement du temps d'arrivée donne
\[
v_+^{\rm arrivée} = 1{,}0036,
\qquad
R^2_{\rm arrivée} = 0{,}99993.
\]
L'ajustement du temps de pic donne
\[
v_+^{\rm pic} = 1{,}0036,
\qquad
R^2_{\rm pic} = 0{,}99980.
\]

Pour la polarisation croix, l'ajustement du temps d'arrivée donne
\[
v_\times^{\rm arrivée} = 1{,}0036,
\qquad
R^2_{\rm arrivée} = 0{,}99980.
\]
L'ajustement du temps de pic donne
\[
v_\times^{\rm pic} = 0{,}9956,
\qquad
R^2_{\rm pic} = 0{,}99974.
\]

Ainsi les deux polarisations se propagent à la même vitesse de graphe,
\[
v_+ \simeq v_\times \simeq 1.
\]

Après calibration des unités de graphe, cela correspond à
\[
v_+ = v_\times = c.
\]

\section{Modes hessiens de l'action BuP}

Pour aller au-delà des perturbations cinématiques, nous linéarisons une action
BuP effective autour d'un graphe d'équilibre \(W_{ij}^{(0)}\). Nous utilisons la
déformation à champ de nœuds contrôlée
\[
W_{ij}(\phi) = W_{ij}^{(0)} \exp\left(\frac{\phi_i + \phi_j}{2}\right).
\]

Le hessien est
\[
H_{ab} = \left. \frac{\partial^2 S_{\rm BuP}}{\partial \phi_a \partial \phi_b} \right|_{\phi = 0}.
\]

Les modes propres satisfont
\[
H \phi_n = \omega_n^2 \phi_n.
\]

Pour un graphe avec \(N = 121\) nœuds, le spectre hessien est strictement positif :
\[
\lambda_{\min} = 0{,}0475,
\qquad
\lambda_{\max} = 10{,}40,
\]
sans modes négatifs. Le graphe de fond est donc linéairement stable.

Les modes bas obéissent à une relation de dispersion effective,
\[
\omega^2 \simeq c_{\rm graphe}^2 \, k^2 + m_{\rm eff}^2.
\]

Pour \(\eta_{\rm lisse} = 1{,}0\), l'ajustement donne
\[
c_{\rm graphe}^2 = 2{,}0289,
\qquad
m_{\rm eff}^2 = 0{,}0352,
\qquad
R^2 = 0{,}99962.
\]
Ainsi
\[
c_{\rm graphe} = 1{,}424.
\]

Un balayage du coefficient de rigidité \(\eta_{\rm lisse}\) donne :

\[
\begin{array}{c|c|c|c}
\eta_{\rm lisse} & \text{modes négatifs} & R^2_{\rm disp} & c_{\rm graphe} \\
\hline
0{,}25 & 0 & 0{,}9948 & 0{,}728 \\
0{,}35 & 0 & 0{,}9972 & 0{,}854 \\
0{,}50 & 0 & 0{,}9985 & 1{,}015 \\
0{,}75 & 0 & 0{,}9993 & 1{,}237 \\
1{,}00 & 0 & 0{,}9996 & 1{,}424 \\
1{,}25 & 0 & 0{,}9998 & 1{,}590 \\
\end{array}
\]

La calibration naturelle est donc proche de
\[
\eta_{\rm lisse} \simeq 0{,}5,
\]
pour laquelle
\[
c_{\rm graphe} \simeq 1.
\]

\section{Tenseur d'émergence et secteur dynamique}

Le régime einsteinien effectif est retrouvé lorsque le tenseur d'émergence BuP
s'annule,
\[
\mathcal H_{\mu\nu}^{\rm BuP} = 0.
\]

Au-delà de la limite lisse, des secteurs supplémentaires peuvent devenir actifs :
\[
\mathcal H_{\mu\nu}^{\rm BuP}
=
\mathcal H_{\mu\nu}^{\rm spec}
+
\mathcal H_{\mu\nu}^{\rm dim}
+
\mathcal H_{\mu\nu}^{\rm nonlocal}
+
\mathcal H_{\mu\nu}^{\rm topo}
+
\mathcal H_{\mu\nu}^{\rm fini}
+
\mathcal H_{\mu\nu}^{\rm dyn}.
\]

Le secteur dynamique est associé aux déviations de propagation :
\[
\mathcal H_{\mu\nu}^{\rm dyn} \neq 0
\quad\Longleftrightarrow\quad
c_{\rm graphe} \neq c.
\]

Dans le vide lisse de l'univers actuel, les observations exigent
\[
c_{\rm graphe} \simeq c.
\]
Ainsi, l'univers actuel est interprété comme étant proche du point fixe de
propagation,
\[
\eta_{\rm lisse} = \eta_\ast,
\qquad
c_{\rm graphe} = c.
\]

\section{Implications cosmologiques exploratoires}

Les résultats suggèrent une hypothèse exploratoire : la vitesse de propagation
des ondes gravitationnelles pourrait sonder l'état informationnel de l'univers.
Dans BuP, cette vitesse n'est pas introduite comme une constante primitive.
Elle est associée aux modes collectifs du graphe d'intrication.

Nous paramétrons une phase primordiale fortement intriquée par
\[
c_{\rm GW}(z) = c\,\alpha(z), \qquad \alpha(z) < 1,
\]
avec
\[
\alpha(z) \to 1
\]
dans le régime lisse actuel.

Une transition simple est donnée par
\[
\alpha(z)
=
\alpha_{\rm early}
+
\frac{1-\alpha_{\rm early}}
{1+\left(z/z_{\rm transition}\right)^s}.
\]
À bas redshift, l'univers est proche du point fixe einsteinien :
\[
\alpha(z) \simeq 1, \qquad c_{\rm GW} = c.
\]
À très grand redshift, la phase fortement intriquée donne :
\[
\alpha(z) \simeq \alpha_{\rm early} < 1.
\]

Nous modélisons l'effet sur le fond stochastique primordial par une fonction de
transfert phénoménologique :
\[
\Omega_{\rm GW}^{\rm BuP}(f)
=
\Omega_{\rm GW}^{\rm GR}(f) \, T_\alpha(f) \, T_{\rm mass}(f).
\]
Le facteur \(T_\alpha\) est pris sous la forme
\[
T_\alpha(\alpha)
=
\alpha^{-p_{\rm accum}}
\exp\left[
-\tau_{\rm damp}
\left(
\frac{1}{\alpha}-1
\right)
\right].
\]
Le premier terme représente l'accumulation du temps de propagation dans une
phase d'intrication dense. Le second terme représente un amortissement
informationnel effectif.

Le facteur \(T_{\rm mass}\) représente un éventuel gap géométrique de basse
fréquence :
\[
T_{\rm mass}(f)
=
\frac{1}{1+\left(f_{\rm mass}/f\right)^q}.
\]
Les tests montrent qu'un gap trop élevé, par exemple
\[
f_{\rm mass}=10^{-6}\,{\rm Hz},
\]
supprime fortement la bande PTA. Pour préserver les PTA, il faut donc prendre
\[
f_{\rm mass}=0
\]
ou
\[
f_{\rm mass}\ll10^{-9}\,{\rm Hz}.
\]

Pour le choix représentatif
\[
\alpha_{\rm early}=0{,}35,
\qquad
z_{\rm transition}=10^{13},
\qquad
p_{\rm accum}=1,
\qquad
\tau_{\rm damp}=0{,}35,
\qquad
f_{\rm mass}=0,
\]
le modèle donne
\[
\alpha_{\rm PTA}=1,
\qquad
\alpha_{\rm LISA}\simeq0{,}35.
\]
La bande PTA reste inchangée par rapport au spectre de référence :
\[
\frac{\Omega_{\rm GW}^{\rm BuP}}
{\Omega_{\rm GW}^{\rm GR}}
\simeq
1{,}00,
\]
tandis que la bande LISA est amplifiée :
\[
\frac{\Omega_{\rm GW}^{\rm BuP}}
{\Omega_{\rm GW}^{\rm GR}}
\simeq
1{,}49.
\]

Un scan sur \(\alpha_{\rm early}\) donne :
\[
\begin{array}{c|c|c|c}
\alpha_{\rm early}
&
\alpha_{\rm PTA}
&
\alpha_{\rm LISA}
&
\Omega_{\rm BuP}/\Omega_{\rm GR}\ {\rm dans\ LISA}
\\
\hline
0{,}20 & 1{,}00 & 0{,}20 & 1{,}233 \\
0{,}30 & 1{,}00 & 0{,}30 & 1{,}473 \\
0{,}35 & 1{,}00 & 0{,}35 & 1{,}492 \\
0{,}50 & 1{,}00 & 0{,}50 & 1{,}409 \\
0{,}70 & 1{,}00 & 0{,}70 & 1{,}230 \\
0{,}90 & 1{,}00 & 0{,}90 & 1{,}069 \\
\end{array}
\]

La réponse LISA est donc non monotone. Le maximum apparaît près de
\[
\alpha_{\rm early}\simeq0{,}35.
\]
Ce maximum n'est pas accidentel. En effet, pour
\[
T_\alpha(\alpha)
=
\alpha^{-p}
\exp\left[
-\tau
\left(
\frac{1}{\alpha}-1
\right)
\right],
\]
on a
\[
\ln T_\alpha
=
-p\ln\alpha
-
\tau
\left(
\frac{1}{\alpha}-1
\right).
\]
Le maximum vérifie
\[
\frac{d}{d\alpha}\ln T_\alpha=0,
\]
donc
\[
-\frac{p}{\alpha}
+
\frac{\tau}{\alpha^2}
=
0.
\]
Ainsi
\[
\alpha_\star=\frac{\tau}{p}.
\]
Avec
\[
p=1,
\qquad
\tau=0{,}35,
\]
on obtient
\[
\alpha_\star=0{,}35,
\]
en accord avec le scan numérique.

Ce résultat suggère que la signature LISA optimale n'est pas obtenue pour une
vitesse primordiale arbitrairement faible, mais pour une phase d'intrication
intermédiaire où l'accumulation temporelle et l'amortissement informationnel
s'équilibrent. Dans ce scénario, les PTA restent compatibles avec le régime
lisse actuel, tandis que LISA pourrait sonder la mémoire d'une phase primordiale
fortement intriquée.

Cette prédiction reste phénoménologique. Elle ne remplace pas un solveur complet
des modes tensoriels cosmologiques. Elle fournit cependant un canal de test
falsifiable pour l'hypothèse BuP :
\[
c_{\rm GW}^{\rm early}<c
\quad
\Longrightarrow
\quad
\text{signature spectrale dans la bande LISA}.
\]

\section{Limitations}

Le présent article est une preuve de concept. Plusieurs limitations subsistent.

Premièrement, l'analyse hessienne utilise une réduction à champ de nœuds,
\[
W_{ij}(\phi) = W_{ij}^{(0)} \exp\left(\frac{\phi_i + \phi_j}{2}\right),
\]
plutôt que le hessien complet au niveau des arêtes
\[
\frac{\delta^2 S_{\rm BuP}}{\delta W_{ij} \delta W_{kl}}.
\]

Deuxièmement, les graphes de fond sont idéalisés. Des graphes astrophysiques ou
cosmologiques réalistes nécessitent une construction dynamique à partir d'états
quantiques microscopiques.

Troisièmement, l'interprétation cosmologique de \(c_{\rm GW}(t)\) est exploratoire.
Elle doit être traitée comme un canal de prédiction, pas comme un résultat établi.

En particulier, la relation effective entre fréquence observée, redshift
d'émission et vitesse \(\alpha(z)\) est modélisée phénoménologiquement et devra
être remplacée par une équation tensorielle cosmologique complète.

\section{Conclusion}

Cet article développe une interprétation BuP des ondes gravitationnelles comme
modes informationnels du graphe d'intrication. Les prototypes numériques et
variationnels montrent que les perturbations du graphe peuvent reproduire les
deux polarisations quadripolaires, se propager comme des modes propagatifs,
émerger de sources dynamiques localisées, et apparaître comme des modes propres
hessiens d'une action BuP effective.

La chaîne centrale est
\[
\delta W_{ij}
\;\rightarrow\;
\delta d_{ij}^{\rm ent}
\;\rightarrow\;
h_{\mu\nu}^{\rm eff}.
\]

Ainsi, les ondes gravitationnelles ne sont pas simplement des ondulations d'une
métrique supposée. Dans BuP, elles sont des sondes du substrat
information-géométrique à partir duquel l'espace-temps émerge.

Enfin, une extension cosmologique exploratoire montre qu'une phase primordiale
fortement intriquée, caractérisée par une vitesse effective réduite
\[
c_{\rm GW}^{\rm early}=c\,\alpha_{\rm early},
\]
peut laisser une empreinte dans le fond stochastique d'ondes gravitationnelles.
Pour \(\alpha_{\rm early}=0{,}35\), le modèle laisse la bande PTA inchangée par
rapport au spectre de référence,
\[
\Omega_{\rm BuP}/\Omega_{\rm GR}\simeq1,
\]
tout en produisant une amplification LISA
\[
\Omega_{\rm BuP}/\Omega_{\rm GR}\simeq1{,}49.
\]
Le scan de \(\alpha_{\rm early}\) montre un maximum près de
\[
\alpha_{\rm early}=0{,}35,
\]
en accord avec l'estimation analytique
\[
\alpha_\star=\tau_{\rm damp}/p_{\rm accum}.
\]
Cette prédiction reste exploratoire, mais elle fournit un canal observationnel
clair pour tester l'hypothèse d'une phase primordiale fortement intriquée.

Les travaux futurs étendront ces outils pour résoudre les régimes d'interaction
non linéaires et explorer la cosmologie primordiale du graphe.

\end{document}
